% ─────────────────────────────────────────────────────────────────────────────
%  SHCT SS2026 — Server Room Cooling  |  Bartha & Kellenberger
% ─────────────────────────────────────────────────────────────────────────────
\documentclass[12pt, a4paper]{article}

% ── Packages ──────────────────────────────────────────────────────────────────
\usepackage[T1]{fontenc}
\usepackage[utf8]{inputenc}
\usepackage[english]{babel}
\usepackage{lmodern}
\usepackage{microtype}
\usepackage{geometry}
\geometry{top=2.5cm, bottom=2.5cm, left=2.5cm, right=2.5cm}

\usepackage{amsmath, amssymb, mathtools}
\usepackage{siunitx}
\sisetup{per-mode=symbol, detect-weight=true}
\usepackage{booktabs}
\usepackage{array}
\usepackage{graphicx}
\usepackage{xcolor}
\usepackage{fancyhdr}
\usepackage{titlesec}
\usepackage{hyperref}
\usepackage{cleveref}
\usepackage{float}
\usepackage[style=numeric, sorting=none, backend=biber]{biblatex}
\addbibresource{references.bib}

% ── Colours ───────────────────────────────────────────────────────────────────
\definecolor{epseblue}{RGB}{0,83,156}

% ── Section colours ───────────────────────────────────────────────────────────
\titleformat{\section}{\large\bfseries\color{epseblue}}{}{0em}{\thesection\quad}
\titleformat{\subsection}{\normalsize\bfseries\color{epseblue}}{}{0em}{\thesubsection\quad}
\titleformat{\subsubsection}{\normalsize\bfseries}{}{0em}{\thesubsubsection\quad}

% ── Header / Footer ───────────────────────────────────────────────────────────
\pagestyle{fancy}
\fancyhf{}
\fancyhead[L]{\small\color{epseblue}\textbf{SHCT SS2026} --- Server Room Cooling}
\fancyhead[R]{\includegraphics[height=1cm]{epse_logo.png}}
\fancyfoot[C]{\thepage}
\setlength{\headheight}{30pt}

% ── Math macros ───────────────────────────────────────────────────────────────
\newcommand{\mdot}{\dot{m}}
\newcommand{\Qdot}{\dot{Q}}
\newcommand{\Wdot}{\dot{W}}
\newcommand{\hstar}{h^{*}}
\newcommand{\Tco}{T_{\mathrm{co}}}
\newcommand{\Tev}{T_{\mathrm{ev}}}
\newcommand{\Tamb}{T_{\mathrm{amb}}}
\newcommand{\Troom}{T_{\mathrm{room}}}
\newcommand{\TAC}{T_{\mathrm{AC}}}
\newcommand{\dTpin}[1]{\Delta T_{\mathrm{pinch,#1}}}
\newcommand{\dTapp}{\Delta T_{\mathrm{app,min}}}
\newcommand{\dTsh}{\Delta T_{\mathrm{sh}}}
\newcommand{\dTsc}{\Delta T_{\mathrm{sc}}}

% ─────────────────────────────────────────────────────────────────────────────
\begin{document}

% ── Title page ────────────────────────────────────────────────────────────────
\begin{titlepage}
\centering
\vspace*{2cm}
{\huge\bfseries\color{epseblue} Simulation of a Server Room Cooling System\par}
\vspace{1cm}
{\Large Sustainable Heat and Cold Technologies (SHCT) --- SS2026\par}
\vspace{2cm}
{\large
  Josua Bartha \quad (20-933-354) \\[0.4em]
  Max Kellenberger \quad (20-961-231)
\par}
\vspace{1cm}
{\large Chair of Energy Systems and Process Engineering (EPSE)\\
ETH Zürich\par}
\vfill
{\large \today\par}
\end{titlepage}

% ─────────────────────────────────────────────────────────────────────────────
\section{Abstract}
\label{sec:abstract}

% TODO: max 200 words
\textit{[Placeholder: summarise motivation, methodology (ODE model, CoolProp performance maps), key results (best refrigerant/diameter, annual energy cost), and main conclusions in at most 200 words.]}

% ─────────────────────────────────────────────────────────────────────────────
\section{Methodology}
\label{sec:method}

% ─────────────────────────────────────────────────────
\subsection{Room Model and Control Strategy}
\label{sec:room}

\subsubsection{Moist-Air State Equations}

The server room (\SI{3}{m} $\times$ \SI{6}{m} $\times$ \SI{10}{m}, $V = \SI{180}{m^3}$) is modelled as a
well-mixed control volume containing a mass $m_{\mathrm{air}}$ of moist air.
Two thermodynamic states are tracked: the specific enthalpy of moist air
$\hstar = h^*(\bar{T},X)$ [\si{kJ/kg_{dry}}] and the mass of dissolved water
$m_w$ [\si{kg}].
Their time evolution follows from the first law and the water balance:

\begin{align}
  m_{\mathrm{air}}\,\frac{d\hstar}{dt}
    &= q_{\mathrm{server}}(t) - \Qdot_{\mathrm{delivered}}(t)
       + \mdot_{\mathrm{vent}}\!\bigl[\hstar_{\mathrm{amb}}(t) - \hstar(t)\bigr],
  \label{eq:hstar_ode}\\[4pt]
  \frac{dm_w}{dt}
    &= \mdot_{\mathrm{vent}}\!\bigl[X_{\mathrm{amb}} - X(t)\bigr].
  \label{eq:mw_ode}
\end{align}

Here $q_{\mathrm{server}}(t)$ [\si{kW}] is the time-varying server heat load,
$\Qdot_{\mathrm{delivered}}$ is the cooling power delivered to the room (from ventilation
or the AC unit), $\mdot_{\mathrm{vent}}$ is the ventilation mass flow rate of dry
air, $X_{\mathrm{amb}}$ is the outdoor specific humidity and $\hstar_{\mathrm{amb}}$
the outdoor moist-air enthalpy.

\subsubsection{Fan Capacity}

A single axial fan (\SI{40}{cm} $\times$ \SI{40}{cm} cross section, face velocity
$v_{\mathrm{fan}} = \SI{5.5}{m/s}$, air density $\rho_{\mathrm{air}} = \SI{1.2}{kg/m^3}$) limits the
supply mass flow rate.
Its maximum capacity is

\begin{equation}
  \mdot_{\mathrm{vent,max}}
  = \rho_{\mathrm{air}}\,v_{\mathrm{fan}}\,A_{\mathrm{fan}}
  = \SI{1.2}{kg/m^3}\times\SI{5.5}{m/s}\times(\SI{0.4}{m})^2
  = \SI{1.056}{kg/s}.
  \label{eq:mdot_max}
\end{equation}

The fan electric power is

\begin{equation}
  \Wdot_{\mathrm{fan}}
  = \frac{\mdot_{\mathrm{vent}}}{\rho_{\mathrm{air}}}\,\frac{\Delta p_{\mathrm{fan}}}{\eta_{\mathrm{fan}}},
  \quad \Delta p_{\mathrm{fan}} = \SI{100}{Pa},\; \eta_{\mathrm{fan}} = 0.60.
\end{equation}

\subsubsection{Ventilation Mode (Free Cooling)}

Ventilation is feasible whenever outdoor air is cool enough to absorb the server heat load
within the fan's flow-rate limit.
The required mass flow to maintain the room enthalpy balance is

\begin{equation}
  \mdot_{\mathrm{req}}
  = \frac{q_{\mathrm{server}}}{\hstar_{\mathrm{room}} - \hstar_{\mathrm{amb}}}.
  \label{eq:mdot_req}
\end{equation}

Ventilation mode is activated if and only if both of the following conditions hold:

\begin{align}
  \Tamb &< \Troom, \label{eq:vent_cond1}\\[2pt]
  \mdot_{\mathrm{req}} &\leq \mdot_{\mathrm{vent,max}}. \label{eq:vent_cond2}
\end{align}

Condition~\eqref{eq:vent_cond1} ensures the outdoor air can absorb heat from the room.
Condition~\eqref{eq:vent_cond2} ensures the required flow is within the fan's physical limit.
If \eqref{eq:vent_cond2} is not satisfied, $\mdot_{\mathrm{vent}}$ is clamped to
$\mdot_{\mathrm{vent,max}}$ and the remaining cooling deficit is covered by the AC compressor.

\subsubsection{AC Thermostat with Hysteresis}

The compressor is governed by a two-point thermostat with asymmetric hysteresis:

\begin{align}
  \Troom \geq T_{\mathrm{ON}}  &= \SI{16}{\degreeCelsius}
    \quad\Rightarrow\quad \text{compressor switched ON},
    \label{eq:ton}\\[2pt]
  \Troom \leq T_{\mathrm{OFF}} &= \SI{10}{\degreeCelsius}
    \quad\Rightarrow\quad \text{compressor switched OFF}.
    \label{eq:toff}
\end{align}

The set-point is $T_{\mathrm{set}} = \SI{15}{\degreeCelsius}$, so the switch-on hysteresis
is $\Delta T_{\mathrm{ON}} = +\SI{1}{K}$ and the switch-off hysteresis is
$\Delta T_{\mathrm{OFF}} = -\SI{5}{K}$.
The wide lower band provides thermal buffer capacity: the room is cooled deeply before the
compressor stops, extending the period until the next switch-on and thereby reducing cycling.

To prevent compressor damage, additional minimum-time constraints are enforced:

\begin{align}
  t_{\mathrm{on}} &\geq \SI{5}{min} \quad \text{(minimum run time after switch-on)},\\[2pt]
  t_{\mathrm{off}} &\geq \SI{10}{min} \quad \text{(minimum standstill time after switch-off)}.
\end{align}

\subsubsection{Overload Detection}

A timestep is flagged as an \emph{overload event} when the compressor is running at full
capacity ($\Qdot_{\mathrm{AC}} = \Qdot_{\mathrm{cool,max}}$), the room temperature is above
the set-point, and the thermal load exceeds the deliverable cooling power:

\begin{equation}
  q_{\mathrm{server}} > \Qdot_{\mathrm{cool,max}}
  \quad\text{and}\quad
  \Troom > T_{\mathrm{overload,min}} = T_{\mathrm{ON}} = \SI{16}{\degreeCelsius}.
  \label{eq:overload}
\end{equation}

% ─────────────────────────────────────────────────────
\subsection{Vapour-Compression Cycle and Performance Maps}
\label{sec:cycle}

\subsubsection{Fixed Design Parameters}

The vapour-compression cycle is described by the fixed parameters in
\Cref{tab:fixed_params}.

\begin{table}[H]
\centering
\caption{Fixed cycle design parameters.}
\label{tab:fixed_params}
\begin{tabular}{llSl}
\toprule
Symbol & Description & {Value} & Unit \\
\midrule
$\TAC$         & Evaporator blow-off temperature (room air leaving evaporator) & 5  & \si{\degreeCelsius} \\
$\dTsh$        & Superheating at compressor inlet                             & 5  & \si{K} \\
$\dTsc$        & Subcooling at condenser exit                                 & 5  & \si{K} \\
$\dTpin{ev}$   & Minimum evaporator approach temperature (pinch)               & 5  & \si{K} \\
$\dTpin{co}$   & Minimum condenser approach temperature (pinch)                & 10 & \si{K} \\
$\Pi_{\min}$   & Minimum pressure ratio $p_{\mathrm{co}}/p_{\mathrm{ev}}$     & 2  & --- \\
\bottomrule
\end{tabular}
\end{table}

\subsubsection{Cycle State Points}

The ideal vapour-compression cycle consists of four state points:
\begin{itemize}
  \item \textbf{State 1} (compressor inlet): superheated vapour at evaporation
        pressure, $T_1 = \Tev + \dTsh$.
  \item \textbf{State 2} (compressor outlet): superheated high-pressure vapour
        at temperature $T_2 > \Tco$ (determined by isentropic compression
        corrected for volumetric efficiency).
  \item \textbf{State 3} (condenser exit): subcooled liquid,
        $T_3 = \Tco - \dTsc$.
  \item \textbf{State 4} (after expansion): two-phase mixture at evaporation
        pressure, $T_4 = \Tev$.
\end{itemize}
All thermodynamic properties ($h$, $p$, $s$, $\rho$, \ldots) are evaluated with
CoolProp~\cite{coolprop}.

\subsubsection{Evaporator Pinch Constraints}

The evaporator transfers heat from the room air (heat source, inlet temperature
$\Troom$, outlet temperature $\TAC$) to the refrigerant (evaporation temperature
$\Tev$, superheated outlet at $T_1$).
Following the approach-temperature formulation for a cross-flow heat exchanger with
an approximately isothermal refrigerant side~\cite{coolprop}, two pinch constraints
must be satisfied:

\begin{align}
  \Troom - T_1
    &\geq \dTapp, \tag{P-ev-1}\label{eq:pev1}\\[2pt]
  \TAC - \Tev
    &\geq \dTapp. \tag{P-ev-2}\label{eq:pev2}
\end{align}

Constraint~\eqref{eq:pev1} ensures that the room air entering the evaporator is warmer
than the superheated refrigerant leaving it; constraint~\eqref{eq:pev2} applies at the
cold end where the air exits at $\TAC$ and the refrigerant enters at $\Tev$.
Since $\TAC$ and $\dTpin{ev} = \dTapp$ are fixed design parameters,
\eqref{eq:pev2} directly fixes the evaporation temperature:

\begin{equation}
  \boxed{\Tev = \TAC - \dTpin{ev} = \SI{5}{\degreeCelsius} - \SI{5}{K} = \SI{0}{\degreeCelsius}.}
  \label{eq:Tev_fixed}
\end{equation}

\eqref{eq:pev1} is automatically satisfied for any room temperature above
$T_1 = \Tev + \dTsh = \SI{5}{\degreeCelsius}$, which is always the case in practice.

\subsubsection{Condenser Pinch Constraints}

The condenser rejects heat from the refrigerant (condensation temperature $\Tco$,
de-superheated inlet at $T_2$, subcooled exit at $T_3$) to the room air acting as
the heat sink at temperature $\Troom$.
Three pinch points arise in a counter-flow condenser~\cite{coolprop}:

\begin{align}
  T_2 - \Troom
    &\geq \dTapp, \tag{P-co-1}\label{eq:pco1}\\[2pt]
  \Tco - T_{\mathrm{sink,dew}}
    &\geq \dTapp, \tag{P-co-2}\label{eq:pco2}\\[2pt]
  T_3 - \Troom
    &\geq \dTapp. \tag{P-co-3}\label{eq:pco3}
\end{align}

Here $T_{\mathrm{sink,dew}}$ is the temperature of the sink-medium at the cross-section
where the refrigerant transitions from superheated to saturated (the ``dew-point
cross-section'' of the condenser).
For an isothermal air-side approximation (cross-flow, well-mixed), all three constraints
reduce to the single binding constraint at the condensation plateau:

\begin{equation}
  \Tco - \Troom \geq \dTpin{co}
  \quad\Longrightarrow\quad
  \Tco \geq \Troom + \dTpin{co}.
  \tag{P-co}\label{eq:pco}
\end{equation}

With $\dTpin{co} = \SI{10}{K}$, the condensation temperature must lie at least
\SI{10}{K} above the current room temperature at every operating point.

\subsubsection{Pressure-Ratio Constraint}

The compressor requires a minimum pressure ratio to operate reliably (valve dynamics,
minimum volumetric efficiency):

\begin{equation}
  \Pi \;=\; \frac{p_{\mathrm{co}}(\Tco)}{p_{\mathrm{ev}}(\Tev)} \;\geq\; \Pi_{\min} = 2.0.
  \tag{C-pr}\label{eq:pressure_ratio}
\end{equation}

Since $\Tev$ is fixed, this constraint defines a lower bound on $\Tco$ that depends on
the refrigerant's vapour-pressure curve $p_{\mathrm{sat}}(T)$.

\subsubsection{Subcritical Constraint}

To maintain a vapour-compression cycle (rather than a transcritical or supercritical
process), the condensation temperature must stay below the refrigerant's critical point:

\begin{equation}
  \Tco < T_{\mathrm{crit}},
  \tag{C-sub}\label{eq:subcrit}
\end{equation}

where $T_{\mathrm{crit}}$ is the critical temperature of the refrigerant
(\SI{96.1}{\degreeCelsius} for R290, \SI{94.7}{\degreeCelsius} for R1234yf,
\SI{127.1}{\degreeCelsius} for DME).

\subsubsection{Condensation-Temperature Optimisation (NLP Formulation)}

For a given operating point $\mathbf{w} = (\Troom, \Tev)$, the COP is maximised by
choosing the lowest feasible $\Tco$.
Following the nonlinear-programming standard form from the pinch-point optimisation
theory~\cite{coolprop}:

\begin{equation}
  \max_{\Tco \in \mathbb{R}}
  \quad \mathrm{COP}(\Tco,\, \mathbf{w})
  \label{eq:nlp_obj}
\end{equation}
\vspace{-1em}
\begin{align}
  \text{subject to}\quad
  g_1(\Tco,\mathbf{w}) &= \Tco - \Troom - \dTpin{co} \;\geq\; 0,
    \tag{NLP-1}\label{eq:nlp1}\\
  g_2(\Tco,\mathbf{w}) &= p_{\mathrm{co}}(\Tco)/p_{\mathrm{ev}}(\Tev) - \Pi_{\min}
    \;\geq\; 0,
    \tag{NLP-2}\label{eq:nlp2}\\
  g_3(\Tco,\mathbf{w}) &= T_{\mathrm{crit}} - \Tco - \epsilon \;\geq\; 0,
    \tag{NLP-3}\label{eq:nlp3}\\[4pt]
  \Tco^{\min} &\leq \Tco \leq \Tco^{\max},
    \tag{NLP-bnd}\label{eq:nlp_bnd}
\end{align}

with $\Tco^{\min} = \max\!\bigl(\Troom + \dTpin{co},\;
  p_{\mathrm{sat}}^{-1}(\Pi_{\min} \cdot p_{\mathrm{ev}})\bigr)$
and $\Tco^{\max} = T_{\mathrm{crit}} - \epsilon$.
Because COP is strictly decreasing in $\Tco$ (lower condensation pressure $\Rightarrow$
smaller enthalpy rise across the compressor), the optimum is always attained at the
lower bound:

\begin{equation}
  \Tco^* = \Tco^{\min}
          = \max\!\Bigl(
              \underbrace{\Troom + \dTpin{co}}_{\text{pinch}},\;
              \underbrace{p_{\mathrm{sat}}^{-1}\!\bigl(\Pi_{\min}\cdot p_{\mathrm{ev}}(\Tev)\bigr)}_{\text{pressure ratio}}
            \Bigr).
  \label{eq:Tco_opt}
\end{equation}

Since $\Tev$ is fixed~\eqref{eq:Tev_fixed}, only $\Troom$ varies at runtime, and the
performance maps are precomputed on a grid of
$\Troom \in [10,30]\,\si{\degreeCelsius}$ and
$\Tamb \in [1,35]\,\si{\degreeCelsius}$
using \eqref{eq:Tco_opt}.

\subsubsection{Compressor Model}

The compressor is a reciprocating piston machine characterised by its piston bore
diameter $D$ (three options: \SI{30}{mm}, \SI{40}{mm}, \SI{50}{mm}) and the
compressor type \texttt{CBar} (clearance-volume volumetric efficiency model).
The refrigerant mass flow through the compressor is

\begin{equation}
  \mdot_{\mathrm{ref}}
  = n_{\mathrm{comp}}\,V_{\mathrm{swept}}\,\rho_1\,\eta_{\mathrm{vol}}(\Pi),
  \label{eq:mdot_ref}
\end{equation}

where $n_{\mathrm{comp}}$ is the rotational speed (fixed), $V_{\mathrm{swept}} \propto D^3$
is the swept volume, $\rho_1 = \rho(\Tev + \dTsh, p_{\mathrm{ev}})$ is the
refrigerant density at the compressor inlet, and $\eta_{\mathrm{vol}}$ is the
volumetric efficiency from the clearance-volume model.
The cycle COP and evaporator capacity follow directly:

\begin{align}
  \mathrm{COP}(\Tco)
    &= \frac{h_1 - h_4}{h_2 - h_1},
    \label{eq:cop}\\[4pt]
  \Qdot_{\mathrm{cycle}}
    &= \mdot_{\mathrm{ref}}\,(h_1 - h_4).
    \label{eq:Qcycle}
\end{align}

\subsubsection{Effective Cooling Capacity with Fan Limit}

The AC unit requires a minimum ventilation flow to remove the condenser heat from the
room.
The maximum evaporator duty is therefore limited by both the refrigerant cycle capacity
and the fan:

\begin{equation}
  \Qdot_{\mathrm{cool,max}}
  = \min\!\Bigl(
      \Qdot_{\mathrm{cycle}},\;
      \mdot_{\mathrm{vent,max}}\,
      \bigl[\hstar_{\mathrm{room}} - \hstar_{\mathrm{AC}}\bigr]
    \Bigr),
  \label{eq:Qcoolmax}
\end{equation}

where $\hstar_{\mathrm{AC}} = \hstar(\TAC, X_{\mathrm{AC,sat}})$ is the enthalpy of
saturated air at the evaporator blow-off temperature~$\TAC$.
The fan-flow constraint thus applies to both operating modes: it directly limits
$\mdot_{\mathrm{req}}$ in ventilation mode~\eqref{eq:vent_cond2} and indirectly caps
$\Qdot_{\mathrm{cool,max}}$ in AC mode via~\eqref{eq:Qcoolmax}.

\subsubsection{Part-Load (On/Off) COP Degradation}

When the server heat load is below $\Qdot_{\mathrm{cool,max}}$, the compressor cycles
on and off.
The part-load ratio is

\begin{equation}
  r = \frac{q_{\mathrm{server}}}{\Qdot_{\mathrm{cool,max}}}, \qquad 0 < r \leq 1.
  \label{eq:plr}
\end{equation}

The effective (result) COP accounting for on/off degradation is

\begin{equation}
  \mathrm{COP}_{\mathrm{result}}
  = \mathrm{COP}_{\mathrm{inner}} \cdot \frac{r}{0.9\,r + 0.1},
  \label{eq:cop_degradation}
\end{equation}

where $\mathrm{COP}_{\mathrm{inner}}$ is the steady-state COP from~\eqref{eq:cop}
and the factor $r/(0.9r+0.1)$ models the parasitic losses during start-up transients
(at $r=1$: $\mathrm{COP}_{\mathrm{result}} = \mathrm{COP}_{\mathrm{inner}}$;
at $r \to 0$: $\mathrm{COP}_{\mathrm{result}} \to 0$).

% ─────────────────────────────────────────────────────
\subsection{Weather Data and Seasonal Weighting}
\label{sec:weather}

% TODO: describe the 5-min weather profiles (source, location, seasons)
\textit{[Placeholder: describe daily ambient temperature and humidity profiles (seasons:
winter 91 d, spring 91 d, summer 91 d, fall 92 d), server load profile source.]}

% ─────────────────────────────────────────────────────
\subsection{Performance Maps and Simulation Procedure}
\label{sec:maps}

% TODO: describe precomputation grid, bilinear interpolation, simulation loop
\textit{[Placeholder: describe COP/\,$\Qdot_{\mathrm{cool,max}}$ maps on
($\Tamb$, $\Troom$) grid (9×6 points), bilinear interpolation at runtime,
5-minute timestep ODE integration.]}

% ─────────────────────────────────────────────────────
\subsection{Annual Energy and Cost}
\label{sec:energy}

The annual compressor energy consumption is accumulated over all four seasons,
each spanning $d_s$ days:

\begin{equation}
  E_{\mathrm{year}}
  = \sum_{s \in \mathcal{S}} d_s
    \sum_{i=1}^{N_{\mathrm{steps}}}
    \Wdot_{\mathrm{comp},i}^{(s)}\,\Delta t,
  \label{eq:Eyear}
\end{equation}

where $\Wdot_{\mathrm{comp},i}^{(s)} = \Qdot_{\mathrm{delivered},i}/\mathrm{COP}_{\mathrm{result},i}$
is the instantaneous compressor power in season $s$ at timestep $i$, and
$\Delta t = \SI{300}{s}$ is the integration timestep.
The annual electricity cost is

\begin{equation}
  K_{\mathrm{year}} = c_{\mathrm{el}}\,E_{\mathrm{year}},
  \qquad c_{\mathrm{el}} = \SI{0.25}{CHF/kWh}.
  \label{eq:Kyear}
\end{equation}

% ─────────────────────────────────────────────────────────────────────────────
\section{Uncertainties}
\label{sec:uncertainty}

% TODO
\textit{[Placeholder: discuss main sources of modelling uncertainty: idealised room
(no thermal mass, no humidity condensation on walls), fixed volumetric efficiency
curve, uncertainty in weather profiles, assumed electricity price.]}

% ─────────────────────────────────────────────────────────────────────────────
\section{References}
\label{sec:refs}

\printbibliography[heading=none]

% ─────────────────────────────────────────────────────────────────────────────
\section{Python Module Descriptions}
\label{sec:modules}

\begin{description}
  \item[\texttt{config.py}]
    Central configuration file. Defines all physical constants, design parameters,
    season weights, cost data, and the auto-generated run tag.
  \item[\texttt{performance\_map.py}]
    Precomputes COP and $\Qdot_{\mathrm{cool,max}}$ on the ($\Tamb$,$\Troom$) grid
    using CoolProp for all refrigerant/diameter combinations.
    Solves the NLP~\eqref{eq:nlp_obj}--\eqref{eq:nlp_bnd} for $\Tco^*$ at each grid
    point.
  \item[\texttt{simulation.py}]
    Integrates the ODE system~\eqref{eq:hstar_ode}--\eqref{eq:mw_ode} with a fixed
    5-minute timestep using the precomputed performance maps and the control logic of
    \Cref{sec:room}.
  \item[\texttt{ranking.py}]
    Evaluates feasibility (no overload events) and ranks all configurations by annual
    cost $K_{\mathrm{year}}$.
  \item[\texttt{selection.py}]
    Compares rankings across different configuration sets to identify which
    refrigerant/diameter combinations change rank.
  \item[\texttt{plots.py}]
    Generates the 4-panel time-series diagnostic plots for each simulation run.
  \item[\texttt{main.py}]
    Top-level entry point: runs performance-map precomputation, simulation, ranking,
    and plotting in sequence.
\end{description}

\end{document}
